\documentclass[11pt,a4paper]{report}
\usepackage[T1]{fontenc}
\usepackage[utf8]{inputenc}
\usepackage[spanish,es-tabla]{babel}
\usepackage{amsmath}
\usepackage{booktabs}
\usepackage{csquotes}
\usepackage[margin=2.8cm]{geometry}
\usepackage[backend=biber,style=apa]{biblatex}
\addbibresource{../docs/references/finance.bib}
\addbibresource{../docs/references/neural.bib}
\addbibresource{../docs/references/books.bib}
\addbibresource{../docs/references/brain.bib}
\addbibresource{../docs/references/macro.bib}
\addbibresource{../docs/references/systems.bib}
\addbibresource{../docs/references/frontier.bib}
\addbibresource{../docs/references/deepseek.bib}
\addbibresource{../docs/references/adversarial.bib}
\usepackage[hidelinks]{hyperref}

\title{MARS-TITAN\\\large Memoria causal adaptativa para predicción bursátil multimodal}
\author{Gonzalo García Lama}
\date{Septiembre de 2026}

\begin{document}
\maketitle
\begin{abstract}
MARS-TITAN plantea una comparación experimental de modelos para predecir
retornos residuales de activos financieros a partir de información multimodal.
La pregunta es si una memoria adaptativa de eventos aporta una mejora frente
a referencias más sencillas cuando se controla la disponibilidad temporal de
los datos y el presupuesto computacional. Se propone trabajar con un subconjunto
de FinMultiTime, validación temporal con avance de ventana, ablaciones y análisis
de incertidumbre.
Este documento es una base de redacción. La preparación de datos y los ensayos
de referencia ya tienen informes separados. La arquitectura MARS-TITAN y la
comparación confirmatoria siguen pendientes.
\end{abstract}
\tableofcontents

\chapter{Introducción y objetivos}
El problema de investigación combina predicción financiera, información
heterogénea y adaptación a cambios de mercado. El valor del estudio reside en
determinar qué componentes aportan información útil y bajo qué condiciones
dejan de hacerlo. El estudio plantea una comparación de soluciones.

Los seis objetivos son preparar datos multimodales con disponibilidad verificable,
definir retornos residuales, diseñar una memoria adaptativa con incertidumbre,
compararla con referencias y ablaciones, evaluar predicciones y simulaciones
financieras y analizar críticamente resultados y restricciones. El protocolo y
sus criterios de aceptación se conservan en la documentación del repositorio.

\chapter{Estado del arte}
FinMultiTime constituye el punto de partida de los datos
\parencite{xu2025finmultitime}. Su alineación por periodos no sustituye una auditoría
de disponibilidad. Titans y los mecanismos de aprendizaje durante inferencia
aportan referencias para la memoria \parencite{behrouz2025titans,sun2025ttt}.
La comparación requiere asimismo referencias sencillas y un control del proceso
de selección \parencite{zeng2023dlinear,bailey2017pbo}.

La revisión detallada se organiza en finanzas, memoria, incertidumbre y fundamentos.
Antes de cerrar este capítulo se integrarán las lecturas con una discusión que
relacione los métodos con la pregunta concreta, sin confundir resultados publicados
con resultados reproducidos.

\chapter{Datos y metodología}
Cada muestra relacionará un activo con un instante de decisión, las modalidades
disponibles y una etiqueta que solo podrá utilizarse después de madurar. La
predicción se construirá con historia conocida y se evaluará mediante ventanas
cronológicas. Ajuste, selección, calibración y evaluación tendrán tramos separados.

Para la convención principal de apertura a cierre de la siguiente sesión,
el retorno residual se define como
\begin{equation}
y_{i,t} = r^{OC}_{i,t+1} - \hat\alpha_{i,t}
          - \hat\beta_{i,t} r^{OC}_{m,t+1}.
\end{equation}
Los coeficientes se estimarán con retornos ya conocidos en el instante de decisión.
El retorno futuro del mercado pertenece a la etiqueta, nunca a las entradas.
La residualización no identifica por sí misma un efecto causal económico.

\chapter{Diseño de la comparativa}
La comparación incluirá referencias ingenuas, regresión regularizada, árboles
potenciados y una red temporal compacta. El modelo propuesto se contrastará con
variantes sin
memoria, con escritura uniforme y sin los componentes adicionales de sorpresa.
Los datos, particiones y presupuestos se mantendrán comparables. La memoria se
reiniciará en cada partición temporal y las actualizaciones respetarán la
maduración de etiquetas.

Se medirán error, dirección, ordenación, calibración y recursos. La simulación usará
retornos de activos y un registro de operaciones con costes de apertura y cierre.
El protocolo fijará las comparaciones principales antes de abrir el conjunto final.

\chapter{Resultados y discusión}
Los ensayos de GRU y DLinear están documentados en el informe de variantes del
repositorio. Sus resultados pertenecen a una muestra dirigida de desarrollo,
no a la comparación confirmatoria. Este capítulo integrará predicciones,
métricas y artefactos identificados por configuración y revisión de código.
Se informarán diferencias, incertidumbre, fallos y costes computacionales.
Las figuras conceptuales no se utilizarán como evidencia de rendimiento.

\chapter{Conclusiones y limitaciones}
Las conclusiones quedan pendientes de la ejecución. Cada objetivo tendrá una
respuesta apoyada en evidencia, o una limitación explícita cuando no pueda
resolverse. El análisis distinguirá mejora observada, resultado inconcluso y
ausencia de mejora, y justificará el trabajo futuro a partir de los límites medidos.

\printbibliography
\end{document}
